\documentclass[11pt]{article}
\usepackage{amsmath,amssymb,amsthm,geometry,hyperref}
\geometry{margin=1in}
\newtheorem{theorem}{Theorem}\newtheorem{lemma}[theorem]{Lemma}\theoremstyle{definition}\newtheorem{definition}[theorem]{Definition}
\title{Enstrophy and Weighted-Enstrophy Balance Identities on $\mathbb{T}^3$:\\ a Lean 4 Formalization}
\author{Benjamin Stanley Frohman\\ ORCID 0009-0006-7068-3718}
\date{October 2026}
\begin{document}\maketitle
\begin{abstract}
We formalize in Lean~4 with Mathlib two exact balance identities from the a-posteriori Lyapunov
functional approach to the 3-D incompressible Navier--Stokes vorticity equation on the torus:
the enstrophy identity $\frac{dE}{dt}=P-\nu\int|\nabla\omega|^2$ and the weighted identity for
$Q=\int|\omega|^4\phi$. Both are proved for every periodic differential algebra whose integral
annihilates spatial derivatives and commutes with the time derivative, for trigonometric polynomials on $\mathbb{T}^3$, and for smooth ($C^\infty$) fields on
$\mathbb{R}\times\mathbb{T}^3$. We also formalize the $L^2$ strain identity, the exact evolution of
$S=E-\frac{\kappa}{2}Q$, a weight-along-paths bound, and boundedness of the comparison ODE. The development has $0$ \texttt{sorry} and uses only
the axioms \texttt{propext}, \texttt{Classical.choice}, \texttt{Quot.sound}.
\end{abstract}

\section{Introduction}
\emph{Commentary (adapted from the author's paper).} Classical energy methods for the
Navier--Stokes equations control the kinetic energy but not the enstrophy in three dimensions;
the vortex-stretching term can produce growth that is difficult to absorb. The approach studies
the exact classical equations together with an auxiliary scalar field $\phi$ transported by
$|\omega|^2$, and the two integrals $E=\tfrac12\int|\omega|^2$ and $Q=\int|\omega|^4\phi$. The first
step is to compute $dE/dt$ and $dQ/dt$ exactly. Those two computations are what this paper
formalizes. They are identities: no inequality and no bound on the solution is involved.

\section{Setting}
Let $u,\omega:\mathbb{T}^3\to\mathbb{R}^3$ and $\phi:\mathbb{T}^3\to\mathbb{R}$ depend on $t$, with
\begin{align}
&\partial_k u_k=0, \label{div}\\
&\partial_t\omega_i=-u_k\partial_k\omega_i+\omega_k\partial_k u_i+\nu\,\partial_k\partial_k\omega_i,\label{vort}\\
&\partial_t\phi=-u_k\partial_k\phi+|\omega|^2,\label{wt}
\end{align}
summing over repeated indices. The relation $\omega=\nabla\times u$ is not needed for either identity.

\begin{definition}[Periodic setting, \texttt{PeriodicSetting}]
A commutative $K$-algebra $A$ with $K$-derivations $\partial_1,\partial_2,\partial_3,\partial_t$, a
$K$-module $B$, a $K$-linear map $I:A\to B$ and a map $D_t:B\to B$ such that
$I(\partial_k f)=0$ for all $k,f$ and $I(\partial_t f)=D_t(If)$.
\end{definition}
Section~\ref{smooth} proves both axioms for smooth periodic fields.

\begin{definition}[Trigonometric model, \texttt{Trig.setting}]
$A=R[\mathbb{Z}^3]$ (finite sums $\sum_n f_n e^{2\pi i n\cdot x}$, $f_n\in R$), where $R$ is a
commutative $\mathbb{C}$-algebra with a derivation $D$. Set
$\partial_k(\sum f_n e_n)=\sum 2\pi i n_k f_n e_n$, $\partial_t(\sum f_n e_n)=\sum D(f_n)e_n$,
$I(f)=f_0$, $D_t=D$.
\end{definition}

\begin{lemma}[\texttt{dx}, \texttt{dt}]
$\partial_k$ and $\partial_t$ are derivations of $R[\mathbb{Z}^3]$.
\end{lemma}
\begin{proof}
Both are linear, so the Leibniz rule reduces to monomials. For $\partial_k$:
$\partial_k(a e_n\cdot b e_m)=2\pi i(n_k+m_k)ab\,e_{n+m}
=a e_n\cdot 2\pi i m_k b e_m+b e_m\cdot 2\pi i n_k a e_n$. For $\partial_t$:
$D(ab)e_{n+m}=(aD b+bDa)e_{n+m}$.
\end{proof}
\begin{lemma}[Torus divergence lemma, \texttt{integral\_dx}, \texttt{integral\_dt}]
$I(\partial_k f)=2\pi i\cdot 0\cdot f_0=0$ and $I(\partial_t f)=D(f_0)=D(If)$.
\end{lemma}
Hence the trigonometric model is a periodic setting, and every theorem below holds for it.

\section{Pointwise lemmas}
\begin{lemma}[\texttt{adv\_eq}]
$u_k\partial_k g=\partial_k(u_k g)-(\partial_k u_k)g$. Hence (\texttt{int\_adv\_eq\_zero}) if
(\ref{div}) holds, $I(u_k\partial_k g)=0$.
\end{lemma}
\begin{proof}Leibniz: $\partial_k(u_kg)=u_k\partial_kg+g\partial_ku_k$. Apply $I$ and use
$I\circ\partial_k=0$ and (\ref{div}).\end{proof}
\begin{lemma}[\texttt{wLapW\_eq}]
$\omega_i\partial_k\partial_k\omega_i=\partial_k(\omega_i\partial_k\omega_i)-\partial_k\omega_i\,\partial_k\omega_i$.
\end{lemma}
\begin{proof}Leibniz on $\partial_k(\omega_i\cdot\partial_k\omega_i)$.\end{proof}
\begin{lemma}[\texttt{two\_w\_adv}]
$2\,\omega_i\,u_k\partial_k\omega_i=u_k\partial_k|\omega|^2$.
\end{lemma}
\begin{proof}$\partial_k(\omega_i\omega_i)=2\omega_i\partial_k\omega_i$.\end{proof}
\begin{lemma}[\texttt{stretch\_eq}]
$\omega_i\omega_j\partial_ju_i=\omega\cdot(\omega\cdot\nabla)u$.
\end{lemma}

\section{Identity (1)}
\begin{lemma}[\texttt{dt\_half\_sq}, \texttt{dt\_vsq\_div}]
Under (\ref{div}), (\ref{vort}):
\[\partial_t|\omega|^2=2\,\omega_i\omega_j\partial_ju_i-2\nu\,|\nabla\omega|^2
+\partial_k\big(2\nu\,\omega_i\partial_k\omega_i-u_k|\omega|^2\big).\]
\end{lemma}
\begin{proof}
$\partial_t|\omega|^2=2\omega_i\partial_t\omega_i$. Substitute (\ref{vort}):
$2\omega_i\partial_t\omega_i=-2\omega_iu_k\partial_k\omega_i+2\omega_i\omega_k\partial_ku_i+2\nu\,\omega_i\partial_k\partial_k\omega_i$.
By \texttt{two\_w\_adv} and \texttt{adv\_eq} with (\ref{div}), the first term is
$-\partial_k(u_k|\omega|^2)$. The second is $2\omega_i\omega_j\partial_ju_i$. By \texttt{wLapW\_eq},
the third is $2\nu\,\partial_k(\omega_i\partial_k\omega_i)-2\nu|\nabla\omega|^2$.
\end{proof}
\emph{Commentary.} The transport term vanishes by periodicity and $\nabla\cdot u=0$; the viscous
term integrates by parts to $-\nu\int|\nabla\omega|^2$; the stretching term remains.
\begin{theorem}[Identity (1), \texttt{enstrophy\_identity}, \texttt{enstrophy\_identity\_T3}]
With $E=\frac12\int|\omega|^2$ and $P=\int\omega_i\omega_j\partial_ju_i$,
\[\frac{dE}{dt}=P-\nu\int_{\mathbb{T}^3}|\nabla\omega|^2 .\]
(Formalized as $D_t\,I(|\omega|^2)=2\,I(\omega_i\omega_j\partial_ju_i)-2\nu\,I(|\nabla\omega|^2)$.)
\end{theorem}
\begin{proof}$D_tI(|\omega|^2)=I(\partial_t|\omega|^2)$; apply $I$ to the lemma; the divergence
term has $I=0$.\end{proof}

\section{Identity (2)}
Write $s=|\omega|^2$, so $|\omega|^4\phi=s^2\phi$.
\begin{lemma}[\texttt{dt\_weighted}]
Under (\ref{vort}), (\ref{wt}):
\[\partial_t(s^2\phi)=4s\phi\,\omega_i\omega_j\partial_ju_i+4\nu\,s\phi\,\omega\cdot\Delta\omega+s^3-u_k\partial_k(s^2\phi).\]
\end{lemma}
\begin{proof}
$\partial_t(s^2\phi)=2s\phi\,\partial_ts+s^2\partial_t\phi$. From the proof of (1),
$\partial_ts=2\omega_i\omega_j\partial_ju_i+2\nu\,\omega\cdot\Delta\omega-u_k\partial_ks$, and by (\ref{wt})
$\partial_t\phi=-u_k\partial_k\phi+s$. So
\[\partial_t(s^2\phi)=4s\phi\,\omega_i\omega_j\partial_ju_i+4\nu s\phi\,\omega\cdot\Delta\omega
-\big(2s\phi\,u_k\partial_ks+s^2u_k\partial_k\phi\big)+s^3,\]
and $2s\phi\,u_k\partial_ks+s^2u_k\partial_k\phi=u_k\partial_k(s^2\phi)$ by Leibniz.
\end{proof}
\emph{Commentary (adapted from the author's paper).} The transport contributions cancel exactly.
Using $\omega\cdot(u\cdot\nabla)\omega=\frac12(u\cdot\nabla)|\omega|^2$ they combine into the pure
divergence $-\int u\cdot\nabla(\phi|\omega|^4)$, which vanishes on $\mathbb{T}^3$.
\begin{lemma}[\texttt{dt\_weighted\_div}]
Under (\ref{div}) as well,
$\partial_t(s^2\phi)=4s\phi\,\omega_i\omega_j\partial_ju_i+4\nu s\phi\,\omega\cdot\Delta\omega+s^3-\partial_k(u_ks^2\phi)$.
\end{lemma}
\begin{proof}\texttt{adv\_eq} with $g=s^2\phi$ and (\ref{div}).\end{proof}
\begin{theorem}[Identity (2), \texttt{weighted\_identity}, \texttt{weighted\_identity\_T3}]
With $Q=\int|\omega|^4\phi$,
\[\frac{dQ}{dt}=4\int|\omega|^2\phi\,\omega\cdot(\omega\cdot\nabla)u+4\nu\int|\omega|^2\phi\,\omega\cdot\Delta\omega+\int|\omega|^6 .\]
\end{theorem}
\begin{proof}$D_tI(s^2\phi)=I(\partial_t(s^2\phi))$; apply $I$, the divergence has $I=0$, and use
\texttt{stretch\_eq}.\end{proof}

\section{Smooth periodic fields}\label{smooth}
\begin{definition}[\texttt{Smooth.A}, \texttt{Smooth.setting}]
$A$ is the algebra of $f\in C^\infty(\mathbb{R}\times\mathbb{R}^3,\mathbb{R})$, $f=f(t,a,b,c)$, with
$f(p+e_k)=f(p)$ for the unit spatial vectors $e_1,e_2,e_3$. $\partial_k f(p)=Df(p)e_k$,
$\partial_tf(p)=Df(p)e_t$, $I(f)(t)=\int_0^1\!\int_0^1\!\int_0^1 f(t,a,b,c)\,dc\,db\,da$, $D_t=\frac{d}{dt}$.
\end{definition}
\begin{lemma}[\texttt{Torus.hasDerivAt\_integral01}]
If $F,\partial_sF$ are jointly continuous on $\mathbb{R}\times\mathbb{R}$, then
$\frac{d}{ds}\int_0^1F(s,x)\,dx=\int_0^1\partial_sF(s,x)\,dx$.
\end{lemma}
\begin{proof}Dominated differentiation: $\partial_sF$ is bounded on the compact set
$\bar B(s_0,1)\times[0,1]$.\end{proof}
\begin{lemma}[\texttt{Torus.hasDerivAt\_torusIntegral}; (b)]
If $f,\partial_tf$ are jointly continuous, then $\frac{d}{dt}\int_{\mathbb{T}^3}f=\int_{\mathbb{T}^3}\partial_tf$.
\end{lemma}
\begin{proof}Apply the previous lemma three times; each inner integral is jointly continuous.\end{proof}
\begin{lemma}[\texttt{Torus.torusIntegral\_deriv\_a/b/c}; (a)]
If $f$ is $C^1$ and $1$-periodic in the differentiated variable, then $\int_{\mathbb{T}^3}\partial_kf=0$.
\end{lemma}
\begin{proof}Differentiate the inner integrals in that variable using the first lemma, then apply
$\int_0^1g'=g(1)-g(0)=0$.\end{proof}
\begin{theorem}[\texttt{Smooth.int\_d}, \texttt{Smooth.int\_dt}, \texttt{Smooth.enstrophy\_identity\_smooth}, \texttt{Smooth.weighted\_identity\_smooth}]
The smooth model is a periodic setting. Hence identities (1) and (2) hold for all
$u,\omega,\phi\in A$ satisfying (\ref{div})--(\ref{wt}).
\end{theorem}

\section{Strain, the functional $S$, the weight, and the comparison ODE}
\begin{theorem}[\texttt{Fourier.L2\_strain\_eq}, \texttt{Fourier.L2\_strain\_le}]
Take $u,\omega,S\in L^2(\mathbb{T}^3)$ whose Fourier coefficients are those of a divergence-free
velocity, its curl, and $S=\frac12(\nabla u+\nabla u^T)$. Then
$\sum_{ij}\|S_{ij}\|_2^2=\frac12\sum_i\|\omega_i\|_2^2$, so $\|S\|_2\le\|\omega\|_2$.
\end{theorem}
\begin{proof}Mode by mode, $\sum|\hat S_{ij}|^2=\frac12|\hat\omega|^2$ when $k\cdot\hat u=0$. Then apply Parseval.\end{proof}
\begin{theorem}[\texttt{Lyapunov.dS\_identity\_smooth}]
Let $S=\frac12\int|\omega|^2-\frac{\kappa}{2}\int|\omega|^4\phi$. For smooth periodic solutions of (\ref{div})--(\ref{wt}),
\[S'=\int\omega_i\omega_j\partial_ju_i-\nu\int|\nabla\omega|^2-\tfrac{\kappa}{2}\Big(4\int|\omega|^2\phi\,\omega_i\omega_j\partial_ju_i+4\nu\int|\omega|^2\phi\,\omega\cdot\Delta\omega+\int|\omega|^6\Big).\]
\end{theorem}
\begin{theorem}[\texttt{Lyapunov.le\_integral\_sq\_of\_deriv}]
Let $h,M$ be continuous, $g(0)=0$, and $g'=h\le M^2$ on $[0,t]$. Then $g(t)\le\int_0^tM^2$.
\end{theorem}
Take $g(s)=\phi(s,X(s))$ along a particle path. Then $g'=|\omega|^2(s,X(s))$, so $\phi\le\int_0^tM^2$
along paths. The chain rule step is not formalized.
\begin{theorem}[\texttt{Lyapunov.ode\_bounded}]
Let $\kappa>0$, and let $y'=y^2(C-\kappa y)$ on $[0,T]$. Then $\min(y_0,C/\kappa,0)\le y(t)\le\max(y_0,C/\kappa,0)$.
\end{theorem}
\begin{proof}Comparison with the barrier $B+\varepsilon$, where $B=\max(y_0,C/\kappa,0)$: at $y=B+\varepsilon$ we have $y'<0$.
Then let $\varepsilon\to0$. The lower bound follows in the same way.\end{proof}
\paragraph{Stated, not proved.} \texttt{CalderonZygmundLp}~$p$ is the bound
$\int|\partial_ju_i|^p\le C\int|\omega|^p$; it is classical for $1<p<\infty$ and is not formalized here.
\texttt{WeightedToSupReduction} says a bound on $S$ implies a bound on $\sup|\omega|^2$.
\texttt{PhiSupControl} is a uniform bound on $|\phi|$. These are \texttt{Prop} definitions only, and
no theorem uses them. The $L^\infty$ Calder\'on--Zygmund bound is false and is not stated.

\section{Scope}
Identities (1) and (2) and the results of the last section are proved as stated. The absorption
step that uses an $L^\infty$ Calder\'on--Zygmund bound, the reduction from $S$ to $\sup|\omega|$, and the
bootstrap are not proved. The paper makes no regularity claim.

\section*{Verification}
\texttt{lake build}: success, 0 \texttt{sorry}. \texttt{lake env lean AxiomAudit.lean}: only
\texttt{propext}, \texttt{Classical.choice}, \texttt{Quot.sound}. Toolchain
\texttt{leanprover/lean4:v4.30.0-rc1}, Mathlib \texttt{v4.30.0-rc1}.
\end{document}
